% sec-01: the shape of a closed day.  Figure 10 (plantuml-type activity with a
% fork and a join), Table 16 (the release list).
\section{The Shape of a Closed Day}

A day with no available counterparty forks into two lanes. One produces things
that must be held, because they reach somebody. The other produces things that
are finished the moment they are written, because they reach nobody until asked
for. Both run at once, and both become releasable at the same point.

\begin{appfloat}[!tb]
\begin{appfig}[plantuml-type / activity with one fork and one join]
% A vertical spine with one fork bar and one join bar, and two lanes at a 4.20cm
% separation.  The lanes leave the fork bar at its left and right thirds and
% rejoin at the corresponding thirds of the join bar, so no lane edge crosses
% the spine.
\node[umlinit] (i) at (4.20,0.90) {};
\node[umlbox,text width=46mm] (start) at (4.20,0.10) {Federal holiday. No counterparty can act};
\node[umlbar,minimum width=62mm] (fk) at (4.20,-0.78) {};
\node[umlstatesoft,text width=34mm] (la) at (2.10,-1.58) {Write the four letters};
\node[umlstatesoft,text width=34mm] (lb) at (2.10,-2.63) {Set the release list};
\node[umlstatesoft,text width=34mm] (lc) at (2.10,-3.68) {Queue the orders, entered nowhere};
\node[umlbox,text width=34mm] (ra) at (6.30,-1.58) {Assemble the data room};
\node[umlbox,text width=34mm] (rb) at (6.30,-2.63) {Fix the recognition letter wording};
\node[umlbox,text width=34mm] (rc) at (6.30,-3.68) {Build the diligence question bank};
\node[umlbar,minimum width=62mm] (jn) at (4.20,-4.50) {};
\node[umlkey,text width=46mm] (ap) at (4.20,-5.32) {Single approval: authorize the release list};
\node[umlctrl,text width=46mm] (nx) at (4.20,-6.24) {Next open session};
\node[umlbox,text width=46mm] (rl) at (4.20,-7.16) {Release the queue, enter the orders};
\node[umlfinal] (f) at (4.20,-7.90) {};
\draw[umlarrow] (i) -- (start);
\draw[umlarrow] (start) -- (fk);
\draw[umlarrow] (2.10,-0.78) -- (la);
\draw[umlarrow] (6.30,-0.78) -- (ra);
\draw[umlarrow] (la) -- (lb);
\draw[umlarrow] (lb) -- (lc);
\draw[umlarrow] (ra) -- (rb);
\draw[umlarrow] (rb) -- (rc);
\draw[umlarrow] (lc) -- (2.10,-4.50);
\draw[umlarrow] (rc) -- (6.30,-4.50);
\draw[umlarrow] (jn) -- (ap);
\draw[umlarrow] (ap) -- (nx);
\draw[umlarrow] (nx) -- (rl);
\draw[umlarrow] (rl) -- (f);
\node[umlguard,text width=34mm,anchor=south] at (2.10,-1.15) {Held lane: reaches a counterparty};
\node[umlguard,text width=34mm,anchor=south] at (6.30,-1.15) {Finished lane: reaches nobody yet};
\node[font=\tiny\itshape,text=fundgray,anchor=west,text width=40mm] at (8.30,-2.63)
  {Everything in the held lane carries a HOLD FOR RELEASE line naming the next
  open session as its earliest send or entry time.};
\end{appfig}
\vspace{-0.60cm}
\figcaption{Figure 10. The closed day as one fork into two concurrent lanes, one of which\\
must be held for release, joined at the single approval the day asks for.}
\end{appfloat}

\subsection{The release list}

\begin{apptable}
\begin{tabularx}{\textwidth}{>{\raggedright\arraybackslash}p{4.0cm} >{\raggedright\arraybackslash}p{2.5cm} >{\raggedright\arraybackslash}X}
\headrow \thead{Item} & \thead{Earliest release} & \thead{Why it cannot go today}\\
\midrule
Congressional district letter & Next open session & District offices observe the federal holiday \\
State finance center letter & Next open session & State offices observe the same holiday \\
Regional development letter & Next open session & The same \\
FDA Pre-Request for Designation & Next open session & The office is closed, and this is the letter least worth sending into a backlog \\
Six queued orders & Next open market session & Both exchanges are closed; an order entered now rejects or fills at an unseen open \\
Grants.gov workspace submission & Next open session & Portal support is closed, and this submission's failure mode needs support \\
eRA Commons profile changes & Next open session & The same \\
\end{tabularx}
\end{apptable}
\vspace{-0.60cm}
\tabcap{Table 16. Seven items with their earliest release and the specific\\
reason each one cannot usefully leave this company on a closed day.}

\subsection{What is finished today and needs no release}

The data room, the recognition letter wording, and the diligence question bank.
None of the three reaches anybody until a counterparty asks for it, so there is
nothing to hold and each is complete on the day it is written. That is the
argument for the day's existence: a closed day cannot produce correspondence that
lands well, but it can produce every internal artifact that correspondence later
depends on.
